\section{Conclusion}
\label{sec:conclusion}

Agent memory has transitioned from a niche subproblem of LLM-based agents to a vibrant, rapidly expanding research field in its own right. Our data-driven survey, curating 952 papers through July 2026, builds on the foundational $3\times 3 \times 3$ taxonomy of \citet{DBLP:journals/corr/abs-2512-13564} and extends it with three new orthogonal dimensions---Temporal Dynamics, Modality, and Biological Inspiration---that capture the emerging themes of forgetting curves, multimodal memory, and neurocognitively inspired architectures.

Our contributions are threefold. First, we describe a fully automated, open-source methodology for maintaining a living paper survey, from YAML-based paper curation through CI validation to interactive web deployment. Second, we demonstrate that the 498 papers published since February 2026 cluster around nine cross-cutting themes that motivate our taxonomy extension, and we propose eight representative systems that illustrate the field's architectural diversity. Third, our gap analysis identifies three persistently sparse taxonomy cells, six open problems, and a three-horizon research roadmap spanning concrete short-term actions (unified benchmarks), medium-term investments (compositional memory, formal trade-off models), and transformative long-term goals (multi-year stability, privacy-compliant forgetting, cross-modal transfer).

The extended taxonomy we propose is intended as a living artifact, not a final statement. As new themes emerge---multi-agent memory sharing, hardware-accelerated memory, neuromorphic architectures---we expect the taxonomy to grow. Our methodology is designed to accommodate such evolution. We release all data, code, and analysis under an open license and invite the community to contribute, critique, and extend this framework.

The central message of this survey is that agent memory research is entering a phase of consolidation. The field has generated a rich diversity of architectures, benchmarks, and insights, but it lacks the standardized evaluation, formal theory, and long-horizon validation that would mark it as a mature discipline. The extended taxonomy and research roadmap presented here are a step toward that maturity.
